\section{Conclusion}

Applying an evaluation cascade to 225 model configurations across three genotoxicity 
datasets shows that, for the Ames mutagenicity benchmark studied here, mixing drug-like and 
industrial compounds---which differ sharply in positive rate and data provenance---produces 
a source heterogeneity gap ($\Delta = 0.390$, scaffold CV to LDO) that is considerably 
larger than what comes from choosing between random and scaffold splitting ($\Delta = 0.046$). 
Threshold decomposition attributed only part of this gap to prior shift; a meaningful 
residual remains that recalibration cannot fix. Algorithm differences are small within a 
single source domain (MCC range = 0.063) but become pronounced under source shift, with RF 
performing considerably better than gradient boosting. \textit{In vivo} micronucleus models can rank 
compounds but do not yet provide reliable binary classification with 2D descriptors. Taken 
together, these results call for a source-aware evaluation cascade as a standard reporting 
requirement for mutagenicity QSAR.
